% !TeX root = ../navier-stokes-manuscript.tex
% ============================================================
% 1.6 A Comparison with Euler
% ============================================================
\subsection{A Comparison with Euler}\label{sub:AComparisonWithEuler}

Participation becomes visible at the boundary between Navier--Stokes and Euler.
The two equations share transport, pressure, and incompressibility,
\[
  \partial_tu+(u\cdot\nabla)u+\nabla p=0,
  \qquad
  \nabla\cdot u=0,
\]
while Euler has no viscous term joining tangential motion across neighbouring layers.
Removing one property from the same physical scene gives a clean test of what that property was doing.

Imagine two rigid boards touching along a flat face and sliding past one another.
Their normal velocities agree because neither board passes through the other.
Their tangential velocities may differ abruptly at the interface.
An inviscid fluid can support the corresponding idealized slip as a weak vortex sheet.

Let the interface be
\[
  \Sigma=\{x_3=0\},
\]
and choose two constants \(U_-\neq U_+\).
Define the stationary velocity and pressure by
\[
  u(x)=
  \begin{cases}
    U_-e_1, & x_3<0,\\
    U_+e_1, & x_3>0,
  \end{cases}
  \qquad
  p(x)=p_0.
\]
The jump \( [u]=(U_+-U_-)e_1 \) is tangent to \(\Sigma\), whose unit normal is \(e_3\).
Consequently, the distributional divergence has no sheet contribution because
\[
  [u]\cdot e_3=0.
\]
The momentum equation is equally transparent in divergence form.
Since
\[
  (u\otimes u)e_3=u(u\cdot e_3)=0
\]
on both sides of the interface, the normal momentum flux has no jump, and the constant pressure contributes no distributional force.
The field is a stationary weak solution of the incompressible Euler equations.

The loss of smoothness remains fully visible.
Writing \(\delta_\Sigma\) for surface measure on the interface, the vorticity is
\[
  \nabla\times u
  =
  (U_+-U_-)e_2\,\delta_\Sigma.
\]
The fluid has concentrated all of its shear into a sheet.
Euler permits that sheet because no term in its momentum law requires tangential momentum to diffuse across \(\Sigma\).

Fixed positive viscosity reads the same interface differently.
In distributions,
\[
  \Delta u
  =
  (U_+-U_-)e_1\,\delta'_\Sigma,
\]
where \(\delta'_\Sigma\) is the normal derivative of the sheet measure.
The stationary transport term still vanishes, while \(\nu\Delta u\) does not.
No pressure can quietly absorb this defect, because the curl of \(\nu\Delta u\) is nonzero whereas the curl of every pressure gradient is zero.
Thus the same discontinuous sheet cannot satisfy the fixed-\(\nu\) Navier--Stokes law.

The comparison now has the exact class-membership form
\[
  \begin{gathered}
  \text{tangential velocity jump}
  \Longrightarrow
  \text{weak Euler vortex sheet},\\
  \text{same sheet with }\nu>0
  \Longrightarrow
  \text{failure of fixed-viscosity participation}
  \Longrightarrow
  \operatorname{Exit}(Q).
  \end{gathered}
\]
This chain belongs to the planar slip model.
It does not say that every nonsmooth scenario is an Euler solution, and the infinite planar sheet is not a finite-energy whole-space counterexample to the Millennium problem.
Its job is local and decisive: the model preserves the shared inviscid rules while exposing exactly why a tangential jump lies outside the fixed-viscosity class.

The velocity jump is the easiest failure to see because it occurs on the lowest rung.
A more subtle endpoint may leave \(u\) finite and continuous while a spatial derivative, an acceleration, or a still higher response loses compatibility.
The class test must follow the equation above velocity itself.
